\documentclass[11pt,a4paper]{article}

\usepackage[utf8]{inputenc}
\usepackage[T1]{fontenc}
\usepackage[spanish,es-nodecimaldot]{babel}
\usepackage{lmodern}
\usepackage{microtype}
\usepackage{geometry}
\usepackage{amsmath,amssymb}
\usepackage{booktabs}
\usepackage{array}
\usepackage{xcolor}
\usepackage{listings}
\usepackage{hyperref}
\usepackage{siunitx}
\usepackage{enumitem}

\geometry{margin=2.4cm}
\hypersetup{colorlinks=true,linkcolor=blue!55!black,urlcolor=blue!55!black,citecolor=blue!55!black}
\sisetup{output-decimal-marker={,},group-separator={.},group-minimum-digits=4}

\definecolor{codebg}{RGB}{246,247,249}
\definecolor{codegreen}{RGB}{25,110,55}
\lstset{
  backgroundcolor=\color{codebg},
  basicstyle=\ttfamily\small,
  keywordstyle=\color{blue!65!black}\bfseries,
  commentstyle=\color{codegreen},
  stringstyle=\color{red!55!black},
  frame=single,
  breaklines=true,
  columns=fullflexible,
  showstringspaces=false,
  tabsize=4
}

\title{\textbf{Optimización conjunta del PLL\_SYS y del divisor PIO\\para síntesis de frecuencias en el RP2040}}
\author{Proyecto PixiePico Firmware}
\date{4 de septiembre de 2026}

\begin{document}
\maketitle

\begin{abstract}
Esta nota presenta un método discreto y exhaustivo para seleccionar conjuntamente
\texttt{REFDIV}, \texttt{FBDIV}, \texttt{POSTDIV1}, \texttt{POSTDIV2} y el divisor
fraccional Q16.8 de una máquina de estados PIO del RP2040. El objetivo es minimizar
el error absoluto de una salida cuadrada cuya frecuencia nominal puede variar. Se
considera un cristal de \SI{12}{MHz}, una instrucción PIO para el estado alto y otra
para el estado bajo, y el intervalo experimental
\(\SI{125}{MHz}\le f_{\mathrm{SYS}}\le\SI{250}{MHz}\). Se presentan resultados para
\SI{7.074}{MHz}, \SI{14.074}{MHz} y \SI{28.074}{MHz}.
\end{abstract}

\section{Planteo del problema}

El oscilador utiliza \texttt{PLL\_SYS} como fuente de reloj de una máquina de
estados PIO. El programa PIO ejecuta dos estados por cada período de salida:
uno mantiene la señal alta y el siguiente la mantiene baja. Por lo tanto,
si \(D\) es el divisor de reloj de la PIO,
\begin{equation}
 f_{\mathrm{out}}=\frac{f_{\mathrm{SYS}}}{2D}.
 \label{eq:pout}
\end{equation}

El divisor PIO se representa en formato Q16.8:
\begin{equation}
 D=d_{\mathrm{int}}+\frac{d_{\mathrm{frac}}}{256}
  =\frac{N}{256},
 \qquad
 N=256d_{\mathrm{int}}+d_{\mathrm{frac}},
 \label{eq:pio_div}
\end{equation}
donde \(0\le d_{\mathrm{frac}}\le255\). Sustituyendo
\eqref{eq:pio_div} en \eqref{eq:pout}:
\begin{equation}
 f_{\mathrm{out}}=\frac{128f_{\mathrm{SYS}}}{N}.
 \label{eq:pout_n}
\end{equation}

\section{Modelo del PLL del RP2040}

Para una referencia externa \(f_{\mathrm{XOSC}}\), el PLL genera
\begin{align}
 f_{\mathrm{ref}} & = \frac{f_{\mathrm{XOSC}}}{REFDIV}, \\
 f_{\mathrm{VCO}} & = f_{\mathrm{ref}}FBDIV, \\
 f_{\mathrm{SYS}} & =
 \frac{f_{\mathrm{VCO}}}{POSTDIV1\,POSTDIV2}.
 \label{eq:pll}
\end{align}

La documentación del RP2040 establece las siguientes restricciones del PLL
\cite{rp2040datasheet}:
\begin{align}
 f_{\mathrm{ref}} &\ge \SI{5}{MHz}, \\
 \SI{750}{MHz}\le f_{\mathrm{VCO}}&\le\SI{1600}{MHz}, \\
 16\le FBDIV&\le320, \\
 1\le POSTDIV1,POSTDIV2&\le7.
\end{align}
Se adopta además \(POSTDIV1\ge POSTDIV2\), tal como recomienda la documentación.

Con \(f_{\mathrm{XOSC}}=\SI{12}{MHz}\), la primera restricción implica
\begin{equation}
 \frac{12}{REFDIV}\ge5
 \quad\Longrightarrow\quad
 \boxed{REFDIV\in\{1,2\}}.
\end{equation}
Por consiguiente, fijar \texttt{REFDIV=1} no era una condición necesaria: era
solamente una reducción del espacio de búsqueda. Permitir \texttt{REFDIV=2}
duplica la resolución disponible en la síntesis del VCO y genera soluciones
adicionales válidas.

El intervalo \SIrange{125}{250}{MHz} para \(f_{\mathrm{SYS}}\) es un requisito
experimental del proyecto PixiePico. No debe confundirse con la frecuencia
nominal garantizada por la especificación del RP2040: la zona superior del
intervalo constituye overclocking y debe validarse sobre el hardware concreto.

\section{Función objetivo conjunta}

Combinando \eqref{eq:pout_n} y \eqref{eq:pll} se obtiene
\begin{equation}
 \boxed{
 f_{\mathrm{out}}=
 \frac{128f_{\mathrm{XOSC}}FBDIV}
 {REFDIV\,POSTDIV1\,POSTDIV2\,N}}
 \label{eq:combined}
\end{equation}

Para una frecuencia objetivo \(f_T\), el problema consiste en encontrar la
tupla entera
\begin{equation}
 (REFDIV,FBDIV,POSTDIV1,POSTDIV2,N)
\end{equation}
que minimiza
\begin{equation}
 E=\left|f_{\mathrm{out}}-f_T\right|,
 \label{eq:objective}
\end{equation}
sujeta a todas las restricciones anteriores y a
\begin{equation}
 \SI{125}{MHz}\le f_{\mathrm{SYS}}\le\SI{250}{MHz}.
\end{equation}

No es apropiado optimizar primero el PLL y luego la PIO de manera independiente.
Una frecuencia \(f_{\mathrm{SYS}}\) que no sea la más cercana a un valor elegido
a priori puede producir un error final menor al combinarse con cierto valor
entero \(N\).

\section{Algoritmo de búsqueda exhaustiva}

El espacio de parámetros es discreto y pequeño. Por ello, una enumeración
exhaustiva entrega el óptimo global dentro del modelo, sin depender de semillas,
gradientes ni mínimos locales.

Para cada cuádrupla válida de parámetros PLL se calcula
\begin{equation}
 N^*=\frac{128f_{\mathrm{SYS}}}{f_T}.
\end{equation}
Como el error es monótono a cada lado de \(N^*\), basta evaluar
\(\lfloor N^*\rfloor\) y \(\lceil N^*\rceil\). No es necesario recorrer los
más de dieciséis millones de valores posibles del registro PIO.

Se recomienda usar aritmética racional o enteros suficientemente anchos. El uso
prematuro de \texttt{float} puede cambiar el orden de candidatos con errores muy
próximos.

\begin{lstlisting}[language=Python,caption={Implementación de referencia del optimizador}]
from fractions import Fraction

XOSC = 12_000_000
SYS_MIN = 125_000_000
SYS_MAX = 250_000_000
VCO_MIN = 750_000_000
VCO_MAX = 1_600_000_000

def optimize(target_hz):
    best = None

    for refdiv in range(1, 64):
        ref = Fraction(XOSC, refdiv)
        if ref < 5_000_000:
            continue

        for fbdiv in range(16, 321):
            vco = ref * fbdiv
            if not (VCO_MIN <= vco <= VCO_MAX):
                continue
            if ref > vco / 16:
                continue

            for postdiv1 in range(1, 8):
                for postdiv2 in range(1, postdiv1 + 1):
                    sys = vco / (postdiv1 * postdiv2)
                    if not (SYS_MIN <= sys <= SYS_MAX):
                        continue

                    ideal_n = 128 * sys / target_hz
                    n0 = ideal_n.numerator // ideal_n.denominator

                    for n in {n0, n0 + 1}:
                        if not (256 <= n <= 0xFFFFFF):
                            continue

                        actual = 128 * sys / n
                        signed_error = actual - target_hz

                        # Mayor VCO como desempate: menor jitter.
                        key = (abs(signed_error), -vco)
                        candidate = (key, refdiv, fbdiv,
                                     postdiv1, postdiv2,
                                     n, vco, sys,
                                     actual, signed_error)

                        if best is None or key < best[0]:
                            best = candidate

    return best
\end{lstlisting}

\section{Resultados}

\begin{table}[htbp]
\centering
\small
\caption{Configuraciones óptimas encontradas mediante búsqueda exhaustiva.}
\label{tab:results}
\begin{tabular}{@{}rrrrrrrr@{}}
\toprule
\(f_T\) (Hz) & REF & FB & PD1 & PD2 & \(N\) & \(f_{\mathrm{SYS}}\) (Hz) & Error (Hz) \\
\midrule
7.074.000  & 2 & 263 & 7 & 1 & 4079 & 225.428.571,429 & +2,731762 \\
14.074.000 & 2 & 211 & 3 & 2 & 1919 & 211.000.000,000 & -3,126628 \\
28.074.000 & 2 & 261 & 4 & 3 & 595  & 130.500.000,000 & -50,420168 \\
\bottomrule
\end{tabular}
\end{table}

\subsection{Caso \(f_T=\SI{7.074}{MHz}\)}

La mejor tupla es
\begin{equation}
 \boxed{(REFDIV,FBDIV,POSTDIV1,POSTDIV2,N)=(2,263,7,1,4079)}.
\end{equation}
De ella se obtiene
\begin{align}
 f_{\mathrm{ref}} &= \frac{\SI{12}{MHz}}{2}=\SI{6}{MHz},\\
 f_{\mathrm{VCO}} &= \SI{6}{MHz}\times263=\SI{1578}{MHz},\\
 f_{\mathrm{SYS}} &= \frac{\SI{1578}{MHz}}{7}
 =\SI{225.428571428571}{MHz}.
\end{align}
El divisor PIO es
\begin{equation}
 D=\frac{4079}{256}=15+\frac{239}{256}=15{,}93359375.
\end{equation}
Finalmente,
\begin{align}
 f_{\mathrm{out}} &= \SI{7074002.731762}{Hz},\\
 \Delta f &= \boxed{+\SI{2.731762}{Hz}},\\
 \varepsilon &= \boxed{+0{,}386169\ \mathrm{ppm}}.
\end{align}

Como comprobación, \SI{233}{MHz} también resulta sintetizable cuando se permite
\texttt{REFDIV=2}: \((2,233,3,2)\) produce un VCO de \SI{1398}{MHz} y
\(f_{\mathrm{SYS}}=\SI{233}{MHz}\). Con \(N=4216\), su error es
\(+\SI{3.795066}{Hz}\), ligeramente mayor que el óptimo global.

\subsection{Caso \(f_T=\SI{14.074}{MHz}\)}

La mejor tupla es
\begin{equation}
 \boxed{(REFDIV,FBDIV,POSTDIV1,POSTDIV2,N)=(2,211,3,2,1919)}.
\end{equation}
Los relojes y el divisor PIO son
\begin{align}
 f_{\mathrm{ref}} &= \SI{6}{MHz},\\
 f_{\mathrm{VCO}} &= \SI{1266}{MHz},\\
 f_{\mathrm{SYS}} &= \frac{\SI{1266}{MHz}}{6}=\SI{211}{MHz},\\
 D &= \frac{1919}{256}=7+\frac{127}{256}=7{,}49609375.
\end{align}
La salida resultante es
\begin{align}
 f_{\mathrm{out}} &= \frac{128(211000000)}{1919}
 =\SI{14073996.873372}{Hz},\\
 \Delta f &= \boxed{-\SI{3.126628}{Hz}},\\
 \varepsilon &= \boxed{-0{,}222156\ \mathrm{ppm}}.
\end{align}

\subsection{Caso \(f_T=\SI{28.074}{MHz}\)}

Una de las tuplas óptimas es
\begin{equation}
 \boxed{(REFDIV,FBDIV,POSTDIV1,POSTDIV2,N)=(2,261,4,3,595)}.
\end{equation}
Los valores asociados son
\begin{align}
 f_{\mathrm{ref}} &= \SI{6}{MHz},\\
 f_{\mathrm{VCO}} &= \SI{1566}{MHz},\\
 f_{\mathrm{SYS}} &= \frac{\SI{1566}{MHz}}{12}=\SI{130.5}{MHz},\\
 D &= \frac{595}{256}=2+\frac{83}{256}=2{,}32421875.
\end{align}
La frecuencia obtenida es
\begin{align}
 f_{\mathrm{out}} &= \frac{128(130500000)}{595}
 =\SI{28073949.579832}{Hz},\\
 \Delta f &= \boxed{-\SI{50.420168}{Hz}},\\
 \varepsilon &= \boxed{-1{,}795974\ \mathrm{ppm}}.
\end{align}

Existen varias tuplas con el mismo error para este objetivo. Por ejemplo, el
mismo VCO de \SI{1566}{MHz} con \((POSTDIV1,POSTDIV2,N)=(6,2,595)\),
\((5,2,714)\) o \((7,1,1020)\) conserva la misma razón final. La selección entre
ellas puede realizarse según frecuencia del sistema, consumo y compatibilidad
con el resto del firmware.

\section{Interpretación y limitaciones}

\begin{enumerate}[leftmargin=*]
 \item La búsqueda obtiene el óptimo global \emph{estático} dentro de las
 restricciones declaradas. No garantiza error cero porque todos los divisores
 relevantes son enteros o racionales con granularidad finita.
 \item El valor nominal calculado no incluye tolerancia, deriva térmica ni
 envejecimiento del cristal de \SI{12}{MHz}. En la práctica, varios ppm del
 cristal pueden superar el error matemático residual del algoritmo.
 \item Una frecuencia racional como \(\SI{1578}{MHz}/7\) es físicamente válida.
 Sin embargo, las API del SDK que almacenan frecuencias en hertz enteros pueden
 reportarla redondeada. La configuración del PLL debe realizarse a partir de
 sus parámetros exactos, no reconstruirse desde el entero redondeado.
 \item Cambiar \texttt{PLL\_SYS} durante la ejecución afecta CPU, PIO y los
 periféricos derivados de ese reloj. Deben revisarse temporizadores, baud rates,
 flash y cualquier lógica que suponga una frecuencia fija. El PLL USB es
 independiente, aunque el firmware USB puede verse afectado indirectamente si
 el cambio no se secuencia correctamente.
 \item Para errores menores que el óptimo estático puede alternarse entre dos
 divisores PIO adyacentes en una proporción calculada. Esto reduce el error
 promedio, pero introduce modulación determinista y componentes laterales; por
 ello constituye otro problema de diseño y no debe confundirse con una
 frecuencia estática exacta.
\end{enumerate}

\section{Recomendación de implementación}

El firmware debería conservar los cinco parámetros del candidato ganador, no
solo \(f_{\mathrm{SYS}}\) y \(D\). La estructura de resultado debe incluir como
mínimo \texttt{refdiv}, \texttt{fbdiv}, \texttt{postdiv1}, \texttt{postdiv2},
\texttt{pio\_div\_int}, \texttt{pio\_div\_frac}, frecuencia calculada y error.
Antes de aplicar el resultado, debe verificarse nuevamente cada restricción y
configurar el PLL mediante sus parámetros exactos.

Para el primer ensayo de hardware se recomienda comenzar con
\SI{14.074}{MHz}, porque su solución usa un \(f_{\mathrm{SYS}}\) entero de
\SI{211}{MHz} y ofrece el menor error relativo de los tres casos. Después debe
medirse la salida con un contador referenciado; sin esa medición, optimizar por
debajo de un ppm es precisión matemática, no necesariamente precisión física.

\begin{thebibliography}{9}
\bibitem{rp2040datasheet}
Raspberry Pi Ltd., \emph{RP2040 Datasheet: A microcontroller by Raspberry Pi},
sección 2.18, ``PLL'', versión de febrero de 2025. Disponible en:
\url{https://datasheets.raspberrypi.com/rp2040/rp2040-datasheet.pdf}.

\bibitem{picosdkpll}
Raspberry Pi Ltd., \emph{Pico SDK---hardware\_pll}, implementación de
\texttt{pll\_init}. Disponible en:
\url{https://github.com/raspberrypi/pico-sdk/blob/master/src/rp2_common/hardware_pll/pll.c}.
\end{thebibliography}

\end{document}
